\section{The architecture below Statement~\ref{st:every}}\label{sec:architecture}

\subsection{The last step}

The file \texttt{CharacterBoxes/MainConclusion.lean} (14 lines) proves Statement~\ref{st:every}:
\leanquote{UpperTheory.FinalTheorems.everyCharacter_of_row}
\leanquote{UpperTheory.FinalTheorems.everyCharacter}
\noindent The first theorem takes one hypothesis, a proposition \lean{UpperTheory.GalStar.Row0b}
about Galois representations (\S\ref{sec:galois}), and uses it three times: twice inside a structure
of ``literature statements'' and once to produce a window bound \lean{UpperTheory.Top.VWin}
(\S\ref{sec:window}). The second theorem discharges the hypothesis with a theorem of the tree,
\lean{UpperTheory.Gal.Int.galoisLocalGlobalStrong_recParameter}.

\subsection{How the literature fields are supplied}

The structure has three fields:
\leanquote{UpperTheory.LiteratureTheorems.OpenLiteratureStatements}
\noindent The first and third fields are the same proposition \lean{UpperTheory.One.M.Row}, which
unfolds to \lean{UpperTheory.GalStar.Row0b} with the local Langlands notion \lean{UpperTheory.Zero.LLC.IsRecT}
\stl{src/Development/AutomorphicArithmetic/Foundations/UnitarySphericalFixing/SatakeParameters/RecT.lean}{6};
the second is the local Langlands correspondence \lean{UpperTheory.Zero.LLC.RecTExists}
(\S\ref{sec:llc}). At \leanlocx{src/Development/AutomorphicArithmetic/CharacterBoxes/MainConclusion.lean}{8}
the structure is filled with $\langle h,\ \texttt{recTExists},\ h\rangle$, where $h$ is the
hypothesis of the first theorem and \lean{UpperTheory.Zero.LLC.recTExists} is a theorem of the tree.

The file declaring this structure contains 195 declarations taking
\texttt{O : OpenLiteratureStatements}. The names of 176 end in \texttt{\_of\_literature}, and many
name published results (Howe duality, the Siegel--Weil formula, the Rallis inner product formula,
Casselman--Wallach, Bernstein--Zelevinsky, Chebotarev, Poitou--Tate, and so on). The textual
dependency check (Appendix~\ref{app:method}) finds 13 declarations that reach a field of
\texttt{O}, directly or through another declaration in the file; for instance,
\lean{UpperTheory.LiteratureTheorems.local_langlands_recT_exists_of_literature} returns the field
\texttt{O.localLanglandsExists}. The other declarations cite results proved elsewhere in the tree.
For example, the body of
\leanquote{UpperTheory.LiteratureTheorems.logDerivOne}
\noindent (elided in the quotation) is \lean{GlobalTheory.LogDerOneProof.logDerOne_holds}. The
final construction supplies the first and third fields with the same Galois-representation
declaration and the middle field with \lean{UpperTheory.Zero.LLC.recTExists}.

\subsection{The map}

Figure~\ref{fig:top} shows the top of the proof. Arrows point from a statement to the statements
its proof uses. The left column is the analytic half (\S\ref{sec:analytic}); the right column is the
automorphic half (\S\ref{sec:window}--\ref{sec:automorphic}); \S\ref{sec:connection} explains the
crossings.

\begin{figure}[ht]
\centering
% Dependency diagram of the top levels (names abbreviated; full names and file:line in the text).
\begin{tikzpicture}[
  every node/.style={font=\scriptsize\ttfamily, align=center},
  thm/.style={draw, rounded corners=2pt, inner sep=2.5pt, minimum height=5mm},
  df/.style={draw, dashed, rounded corners=2pt, inner sep=2.5pt, minimum height=5mm},
  ->, >={Stealth[length=1.6mm]}, node distance=4.5mm and 6mm]
\node[thm] (rh) {riemannHypothesis\\(Main.lean)};
\node[thm, below=of rh] (every) {FinalTheorems.everyCharacter\\$\vdash$ EveryCharacter (Statement 3.1)};
\node[thm, below=of every] (lit) {every\_character\_of\_literature\_of\_vWin\\(O : OpenLiteratureStatements) (hV : VWin)};
\node[thm, below=of lit] (theta) {every\_character\_of\_theta\_integrals\_local\_matching};
\node[thm, below=of theta] (imp) {Top.imprimitive $\circ$ zero\_free\_all\_of\_green\_v\_win\\ZeroFreeAll $\Rightarrow$ EveryCharacter};
\node[thm, below left=6mm and -14mm of imp] (topg) {zeroFreeAll\_of\_topG\\(la = 3/2)};
\node[thm, below left=5mm and -6mm of topg] (landau) {zeroFreeOn\_of\_tiedWin\_cyc\\Landau criterion (formB)\\Chebyshev on windows};
\node[df, below right=5mm and -8mm of topg] (topgdef) {TopG: family classes\\connected};
\node[thm, below right=6mm and -10mm of imp] (supply) {supplyChi\\VWinAt $\psi$ 401 $\Rightarrow$ HostExW (3/2)};
\node[thm, below=of topgdef] (green) {shape\_closed\_link\_n\_of\_green\\(Green's cuspidal characters)};
\node[thm, below=13mm of supply] (vwin) {One.vwin\_of\_row\\$\vdash$ VWin (Sha-length bounds)};
\node[thm, right=8mm of lit] (row) {galoisLocalGlobalStrong\_recParameter\\$\vdash$ Row0b recT};
\node[thm, below=of row] (galth) {galoisRepresentations\_of\_thetaLift};
\node[thm, below left=5mm and -18mm of galth] (tlift) {TP.K.tLiftK\\(theta lifts)};
\node[thm, below right=5mm and -14mm of galth] (geo) {Geo.geoOne\_holds};
\node[thm, right=8mm of every] (llc) {Zero.LLC.recTExists\\(local Langlands)};
\draw (rh) -- (every);
\draw (every) -- (lit);
\draw (lit) -- (theta);
\draw (theta) -- (imp);
\draw (imp) -- (topg);
\draw (imp) -- (supply);
\draw (topg) -- (landau);
\draw (topg) -- (topgdef);
\draw (topgdef) -- (green);
\draw (supply) -- (vwin);
\draw (every) -- (row);
\draw (every.east) -- (llc);
\draw (row) -- (galth);
\draw (galth) -- (tlift);
\draw (galth) -- (geo);
\draw[dashed] (vwin.east) to[out=0, in=-90] (row.south east);
\draw[dashed] (lit.east) -- (row.west);
\end{tikzpicture}

\caption{The top of the proof of \lean{riemannHypothesis}. Boxes are Lean declarations
(abbreviated); solid boxes are theorems, dashed boxes are definitions. Every node is quoted or cited
with file:line in the text.}
\label{fig:top}
\end{figure}
